% !TeX program = xelatex
% vim: sw=0
\documentclass{resume}

\usepackage{hyperref}

% temporally raise textbox for hangul; balancing with eng name
\name{\hangulparens Yojun Moon \raisebox{0.03em}{(문요준)}}

% \mobile{\href{tel:+821040729190}{010-4072-9190}}
\email{\href{mailto:yojun313@postech.ac.kr}{yojun313@postech.ac.kr}}
\github{\href{https://github.com/yojun313}{github.com/yojun313}}
\begin{document}

\begin{rSection}{Education}
	\begin{rEducation}{\textbf{POSTECH (Pohang University of Science and Technology)}}{Mar 2023 -- Present}
		B.S. in Computer Science and Engineering
	\end{rEducation}
\end{rSection}

\begin{rSection}{Work Experience}
	\begin{rSubsection}{\href{https://www.samsungfire.com}{Samsung Fire \& Marine Insurance}}{Jul 2026 -- Aug 2026}
		\textbf{AI Engineer Intern, AI Development Part}
		\\ [0.2em]
		Built two computer vision modules for insurance claim data, each delivered as an end-to-end pipeline with a web service.

		% \item Automatic License Plate Masking for Dashcam Video:
		% \vspace{-0.5em}
		% \begin{list}{--}{}
		% 	\itemsep -0.5em
		% 	\item Delivered a module that masks every vehicle license plate in dashcam footage, reaching 99.0\% recall and 0.98 F1 on an 11.6K-image evaluation set.
		% 	\item Sustained 27+ fps end-to-end, shipped as a web service covering upload, processing, and download.
		% \end{list}

		% \item Detection of Generative-AI Document Forgery:
		% \vspace{-0.5em}
		% \begin{list}{--}{}
		% 	\itemsep -0.5em
		% 	\item Built FinForge, a 6,000-image forgery dataset of Korean insurance claim forms, and fine-tuned TruFor (CVPR 2023) on it.
		% 	\item Raised forgery localization F1 from 0.51 to 0.97 and document-level classification accuracy from 55\% to 93\%, holding 0.88 F1 on low-quality images.
		% \end{list}
		% \item Skills: PyTorch, Ultralytics YOLO, OpenCV, FFmpeg, FastAPI
	\end{rSubsection}

	\begin{rSubsection}{Korean National Police University -- \href{https://knpu.re.kr/}{FPEI (Future Public Safety Engineering Institute)}}{Nov 2023 -- Present}
		\textbf{Research System Developer}
		\\ [0.2em]
		Core developer of the institute's research platform: 8 services on a shared FastAPI backend, deployed in production at {\small\href{https://knpu.re.kr/systems}{*.knpu.re.kr}} ({\small\href{https://github.com/yojun313/knpu}{github.com/yojun313/knpu}}).

		\item Data Collection --- \textit{CRAWLER}, an asynchronous, queue-based crawler with proxy IP rotation, gathering about 1,500 articles and 60,000+ comments per hour from Naver News, and also supporting Naver Blog, Naver Cafe, and YouTube.
		\item Analysis --- \textit{STATISTICS} (statistical analysis and visualization), \textit{NETWORK} (social network extraction and visualization from natural language), and \textit{KEMKIM} (future signal prediction), unified by \textit{MANAGER}, the research resource and big data platform.
		\item AI Legal Complaint Drafting --- built an LLM-based service that drafts complaints meeting legal requirements; received a Letter of Appreciation from the Korean National Police University (Apr 2024).
		\item Skills: Python, FastAPI, MongoDB, vLLM, PyQt, PM2
	\end{rSubsection}
	\begin{rSubsection}{\href{https://sites.google.com/view/haiv/}{POSTECH HAiV Lab}}{Jan 2026 -- Feb 2026}
		\textbf{Undergraduate Researcher}
		\\ [0.2em]
		Researching ACCIDENTOR, a multi-VLM framework that estimates traffic accident fault ratios from dashcam video.
		\item Skills: Python, PyTorch, Vision-Language Models
	\end{rSubsection}

	\begin{rSubsection}{Startup Team (founded by a POSTECH CSE alumnus)}{Feb 2025 -- Jul 2025}
		\textbf{Backend Developer Intern}
		\\ [0.2em]
		Developed the backend of a social networking service and supported data-driven product decisions.

		\item Key Contributions:
		\vspace{-0.5em}
		\begin{list}{--}{}
			\itemsep -0.5em
			\item Developed the SNS backend with NestJS.
			\item Built web crawlers and analyzed the collected data.
			\item Completed the 2025 I-Corps program, including 3 weeks of entrepreneurship training at UC Berkeley (KIC Tech Frontier Program).
		\end{list}
		\item Skills: NestJS, TypeScript, Python, MongoDB, Crawler
	\end{rSubsection}

\end{rSection}

\begin{rSection}{Projects}
	\begin{rSubsection}{\href{https://pcss.postech.ac.kr/}{PCSS (POSTECH Computer Scientist Search)}}{Jan 2025 -- Mar 2026}

		\textit{In-house service for the POSTECH CSE department.} Faculty candidate search service, developed on commission.

		\item Key Implementations:
		\vspace{-0.5em}
		\begin{list}{--}{}
			\itemsep -0.5em
			\item Parsed DBLP XML dumps and stored normalized publication records in MongoDB.
			\item Identified Korean-name authors at scale with an LLM API.
			\item Reduced the search for Korean authors of top-venue papers within a given year and venue range from hours to seconds (several hundred times faster).
		\end{list}
		\item Deployed at \href{https://pcss.postech.ac.kr/}{pcss.postech.ac.kr}
		\item Skills: Python, MongoDB, LLM APIs
	\end{rSubsection}

	\begin{rSubsection}{\href{https://commeet.postech.ac.kr}{Commeet}}{Aug 2026 -- Present}

		\textit{In-house service for the POSTECH CSE department.} Group scheduling service that collects everyone's availability and assigns the meeting times.

		\item Key Implementations:
		\vspace{-0.5em}
		\begin{list}{--}{}
			\itemsep -0.5em
			\item Automatic time slot assignment --- reduced the conflict-free schedule to a degree-constrained bipartite matching and solved it optimally with min-cost max-flow, maximizing the number of people scheduled while packing them into as few sessions as possible.
			\item Contact management --- a reusable address book for frequent invitees, covering both registered users and people who have not signed up.
			\item Professor-only tools --- consultation records in the four categories required by POVIS, with per-student notes, a drawn signature, and a printable form; plus student guidance fee settlement with receipt upload. All records are encrypted at rest, as they hold student personal data.
		\end{list}
		\item Deployed at \href{https://commeet.postech.ac.kr}{commeet.postech.ac.kr}
		\item Skills: FastAPI, MongoDB (Motor), Jinja2, JavaScript
	\end{rSubsection}

	\begin{rSubsection}{\href{https://popo.poapper.club/}{POPO}}{Jan 2026 -- Present}

		Campus web service operated by PoApper, POSTECH's student developer club.

		\item Developing the backend with NestJS.
		\item Skills: NestJS, TypeScript
	\end{rSubsection}
\end{rSection}

\begin{rSection}{Activities}
	\begin{rSubsection}{POSTECH Broadcasting Station (PBS)}{Mar 2023 -- Dec 2025}
		Member of the Programming \& Production Department.
		\item Filmed, edited, and live-broadcast campus entertainment programs.
	\end{rSubsection}

	\begin{rSubsection}{Programming Tutor}{Jan 2025 -- Present}
		Teaching Python and C to middle and high school students.
		\item Topics include object-oriented programming, MongoDB, and FastAPI.
	\end{rSubsection}
\end{rSection}

\begin{rSection}{Skills}
	\begin{tabular}{@{}p{0.25\linewidth}p{0.7\linewidth}}
		Programming Languages
		 & Python, C, C++, TypeScript, JavaScript                          \\ [0.3em]

		AI / Machine Learning
		 & PyTorch, Ultralytics YOLO, OpenCV, vLLM, LLM APIs               \\ [0.3em]

		Backend
		 & FastAPI, NestJS, MongoDB                                        \\ [0.3em]

		Tools / Miscellaneous
		 & Docker, Git, Linux, PyQt, FFmpeg                                \\ [0.3em]
	\end{tabular}
\end{rSection}

\end{document}
